\begin{tikzpicture}[
    font=\sffamily\small,
    >=Stealth,
    arr/.style   = {->,  draw=gray!70, line width=1.2pt},
    biarr/.style = {<->, draw=gray!70, line width=1.2pt},
    titlebox/.style = {
      rectangle, rounded corners=3pt,
      draw=cBlue, line width=1.6pt, fill=cTitle,
      minimum width=12cm, minimum height=0.85cm,
      text centered, text width=11.8cm
    },
    probbox/.style = {
      rectangle, rounded corners=3pt,
      draw=cBlue!80, line width=1pt, fill=cProb,
      minimum width=3.2cm, minimum height=0.82cm,
      text centered, text width=3.0cm
    },
    methbox/.style = {
      rectangle, rounded corners=3pt,
      draw=cOrange, line width=1pt, fill=cMeth,
      minimum width=2.9cm, minimum height=0.72cm,
      text centered
    },
    labelbox/.style = {
      rectangle, rounded corners=2pt,
      draw=gray!60, line width=0.8pt, fill=cLabel,
      minimum width=1.3cm, minimum height=0.7cm,
      text centered, font=\footnotesize\sffamily
    },
  ]
   
  % ── 研究目标 ────────────────────────────────────────────────
  \node[titlebox] (title) at (0,0)
    {弱可压缩不混溶两相流高精度数值模拟方法研究};
  \node[labelbox, left=5mm of title] {研究\\目标};
   
  % ── 科学问题 ────────────────────────────────────────────────
  \node[probbox, below left=1.0cm and 4.1cm of title]  (p1)
    {非物理压力振荡\\（Abgrall 原则）};
  \node[probbox, below=1.0cm of title]                  (p2)
    {分相密度正性保持\\（双曲性条件）};
  \node[probbox, below right=1.0cm and 4.1cm of title] (p3)
    {数值耗散与\\相界面弥散};
  \node[labelbox, left=5mm of p1] {科学\\问题};
   
  \draw[arr] (title.south) -- (p1.north);
  \draw[arr] (title.south) -- (p2.north);
  \draw[arr] (title.south) -- (p3.north);
   
  % ── 研究方法 ────────────────────────────────────────────────
  \node[methbox, below=0.9cm of p1] (m1) {理论分析};
  \node[methbox, below=0.9cm of p2] (m2) {互相验证};
  \node[methbox, below=0.9cm of p3] (m3) {数值实验};
  \node[labelbox, left=5mm of m1]   {研究\\方法};
   
  \draw[arr]   (p1) -- (m1);
  \draw[arr]   (p2) -- (m2);
  \draw[arr]   (p3) -- (m3);
  \draw[biarr] (m1) -- (m2);
  \draw[biarr] (m2) -- (m3);
   
  \end{tikzpicture}